% LaTeX-Kopf fuer die PDF-Fassung der Technischen Dokumentation.
%
% Wird von technische-doku-erzeugen.sh ueber "pandoc -H" eingebunden. Jede
% Zeile hier loest ein Problem, das beim Setzen tatsaechlich aufgetreten ist -
% wer eine entfernt, bekommt das Problem zurueck.

% --- Bilder bleiben, wo sie im Text stehen ---------------------------------
% Ohne float/[H] verschiebt LaTeX ein Diagramm an den naechsten Seitenanfang,
% an dem es passt. Bei zwoelf Diagrammen wandern sie damit reihenweise ans
% Kapitelende, weg von dem Absatz, der sie erklaert.
\usepackage{float}
\let\origfigure\figure
\let\endorigfigure\endfigure
\renewenvironment{figure}[1][2]{\origfigure[H]}{\endorigfigure}

% --- Eine Querformatseite fuer die Gesamtuebersicht ------------------------
% Die Gesamtuebersicht aller Klassen ist rund 3500 Pixel breit. Auf einer
% Hochformatseite schrumpft sie auf Unlesbarkeit. In der Markdown-Quelle
% steht das Bild deshalb zwischen zwei Rohblocks:
%
%     ```{=latex}
%     \begin{landscape}
%     ```
%
% Wichtig: als abgesetzte ```{=latex}-Bloecke, nicht als blanke
% \begin{landscape}-Zeilen. Pandoc zieht sonst alles zwischen \begin und \end
% als eine einzige Roh-LaTeX-Passage zusammen - die Bildsyntax dazwischen
% bleibt unuebersetzt stehen, und die Unterstriche im Dateinamen brechen den
% Lauf mit "Missing $ inserted".
\usepackage{pdflscape}

% --- Lange Bezeichner brechen statt ueber den Rand zu laufen ---------------
% Inline-Quelltext wie SimulierteKassenGegenstelle ist ein einziges Wort ohne
% Trennstelle. LaTeX kann es nicht umbrechen und schiebt es in den Rand, wo es
% abgeschnitten wird - im Satz fehlte genau so das Ende eines Klassennamens.
% seqsplit erlaubt den Umbruch an jeder Stelle; gebrochen wird trotzdem nur
% dort, wo es die Zeile verlangt.
%
% Das \ttfamily muss mit: \seqsplit allein setzt das Argument in der
% Grundschrift, und der Unterschied zwischen Fliesstext und Quelltext ginge
% verloren.
\usepackage{seqsplit}
\let\oldtexttt\texttt
\renewcommand{\texttt}[1]{{\ttfamily\seqsplit{#1}}}
